% ============================================================================
%  chapter9.tex
%  Process Control: Groups, Sessions, and Daemons
% ============================================================================

\chapter{Process Control and Daemons}

\section{From Single Processes to a Community}

Chapter~5 created processes; Chapter~7 connected them with pipes;
Chapter~6 taught them to react. This chapter asks the next question
up: \emph{how does a collection of processes live together?} Real
machines run hundreds of processes at once, and the operating system
must group them so that it can signal them, schedule them, and police
them as units. The two grouping concepts are the \emph{process group}
and the \emph{session}, and they are the backbone of every shell,
every terminal, and every daemon.

\begin{verbatim}
   The process tree. Every process has a parent; the root is init:

     init (pid 1)
       |
       +-- systemd (pid 1's job: supervise everything)
             |
             +-- your login shell (a session leader)
                   |
                   +-- emacs      (foreground process group)
                   +-- make       (a background group)
                   |     +-- gcc
                   |     +-- cc1
                   +-- server     (another background group)

   The tree is one fact; the groups and sessions are the labels the
   kernel uses to decide who gets which signal and who may read the
   terminal.
\end{verbatim}

\begin{osconcept}
The kernel does not manage processes one by one; it manages them as
units. Groups exist so that a signal or a job can address many
processes at once, and sessions exist so that the kernel can tie one
set of processes to one terminal. Deciding what those units are---and
how a process joins or leaves them---is operating system bookkeeping,
and it is what makes job control and daemons possible at all.
\end{osconcept}

\section{Process Groups}

A \emph{process group} is a set of processes that share a group id. The
group's purpose is to be a target: \emph{you can signal a whole group
at once}. When you press Ctrl-C, the terminal does not signal your one
process; it signals your entire foreground \emph{group}---which is why
a pipeline like \code{cat f | sort} stops as one unit.

\begin{lstlisting}[caption={Who am I, and who is my group?}]
#include <stdio.h>
#include <unistd.h>

int main(void) {
    printf("pid       = %d\n", getpid());
    printf("process group = %d\n", getpgrp());
    return 0;
}
\end{lstlisting}

A child inherits its parent's process group; a session leader, or a job
that a shell places in its own group, has a group id equal to its own
pid---it starts a new group alone. The calls that move processes between groups are
\code{setpgid} (join or create a group) and \code{getpgrp} (ask who
your group is). The shell uses \code{setpgid} to gather a pipeline's
children into one group so it can signal them together.

\begin{lstlisting}[caption={Signal a whole group: note the minus sign.}]
#include <stdio.h>
#include <stdlib.h>
#include <unistd.h>
#include <signal.h>

/* Ask every process in "pgrp" to terminate. */
void kill_group(pid_t pgrp) {
    int r;

    r = kill(-pgrp, SIGTERM);    /* -pgrp means "the group named pgrp" */
    if (r < 0) {
        perror("kill group");
    }
}
\end{lstlisting}

\begin{verbatim}
   kill(-pgrp, sig)  vs  kill(pid, sig):

     kill(pid, sig)    signal ONE process
     kill(-pgrp, sig)  signal EVERY process in the group pgrp
     kill(0, sig)      signal every process in MY group
     kill(-1, sig)     signal every process I may signal (extreme)

   This is why Ctrl-C feels like it hits "the whole job": the terminal
   driver runs kill(-foreground_pgrp, SIGINT). One call, many victims.
\end{verbatim}

\begin{osconcept}
Process groups are the kernel's unit of \emph{collective action}. The
OS lets one call address many processes because stopping or starting a
job should be atomic with respect to the whole job, not piecemeal.
The group id is kernel state; \code{setpgid} is a syscall that
reassigns it, and only the kernel can guarantee that a group's members
are tracked consistently.
\end{osconcept}

\section{Sessions and the Controlling Terminal}

A \emph{session} is a collection of process groups, headed by a
\emph{session leader} (the process that created the session). A
session may own one \emph{controlling terminal}. The session is the
unit that owns a terminal; within the session, exactly one process
group---the \emph{foreground group}---is allowed to read it.

\begin{verbatim}
   One session, one terminal, three groups:

      /dev/tty1  (the controlling terminal)
         |
         +---- session led by the login shell (sid = 100)
                |
                +---- FOREGROUND group  (your `cat bigfile`)
                |       -- the ONLY group allowed to read the tty
                |       -- receives Ctrl-C as SIGINT
                |
                +---- background group  (your `sleep 100 &`)
                |       -- NOT allowed to read the tty
                |       -- does NOT receive Ctrl-C
                |
                +---- background group  (your `make &`)

   Two kernel-enforced rules that follow from this picture:

     1. A background group that tries to READ the terminal is STOPPED
        by SIGTTIN. Try it: run `cat` in the background of a shell
        and it freezes, because the kernel refuses to let it steal
        your keystrokes.

     2. When you log out, the terminal closes; the kernel sends SIGHUP
        (hang up) to the session's leader, which the shell uses to
        clean up its jobs. Daemons, as we will see, detach from all
        of this on purpose.
\end{verbatim}

\begin{lstlisting}[caption={Forking a background job is group arithmetic.}]
#include <stdio.h>
#include <stdlib.h>
#include <unistd.h>
#include <sys/types.h>
#include <sys/wait.h>

int main(void) {
    pid_t pid;
    pid_t my_group;

    pid = fork();
    if (pid < 0) {
        perror("fork");
        exit(1);
    }

    if (pid == 0) {
        /* Child: put itself in its own new group, like a shell
           backgrounding a job. Then run a command.                 */
        setpgid(0, 0);            /* new group, id = my pid */
        execlp("sleep", "sleep", "100", NULL);
        perror("execlp");
        exit(127);
    }

    /* Parent: also record the child in its new group, then wait. */
    my_group = getpgrp();
    setpgid(pid, pid);
    printf("parent in group %d, child %d is in its own group %d\n",
           (int) my_group, (int) pid, (int) pid);

    waitpid(pid, NULL, 0);
    return 0;
}
\end{lstlisting}

\begin{verbatim}
   Why both parent AND child call setpgid? A race:

     - if the child runs first, IT calls setpgid(0,0) before the
       parent's setpgid(pid,pid);
     - if the parent runs first, the parent's call fixes the group
       before the child execs.

   Both orders are correct because setpgid is idempotent and allowed
   for a process still in the same session. Shells do exactly this to
   avoid a window where a backgrounded child could read the terminal.
\end{verbatim}

\begin{osconcept}
Sessions and the controlling terminal are the kernel's answer to a
very concrete problem: many processes share one screen, and the screen
must have exactly one reader and one interrupt target at a time. The
kernel owns the terminal device, so it can enforce the foreground rule
mechanically---stopping an intruder with \code{SIGTTIN} and routing
\code{SIGINT} only to the foreground group. This is why job control
feels ``instant and fair'': it is hardware-level policy.
\end{osconcept}

\section{setsid: Leaving the Terminal Behind}

\code{setsid} is the escape hatch. It creates a brand-new session,
makes the caller its leader, and---crucially---\emph{detaches the
caller from any controlling terminal}. After \code{setsid}, your
process is a session leader with no terminal, no foreground group, and
no route for Ctrl-C to reach it. This is the first step of becoming a
daemon.

\begin{lstlisting}[caption={setsid: new session, no terminal.}]
#include <stdio.h>
#include <unistd.h>

int main(void) {
    pid_t sid;

    sid = setsid();
    if (sid < 0) {
        perror("setsid");
        return 1;
    }

    printf("new session id = %d\n", (int) sid);
    printf("(I am now a session leader with no controlling terminal)\n");
    return 0;
}
\end{lstlisting}

The catch: \code{setsid} fails if the caller is already a process
group leader, because a group leader could accidentally make its whole
group leave together. That failure is the entire reason for the
\emph{first fork} in daemonization (next section).

\begin{osconcept}
\code{setsid} is the kernel's ``give me a clean room'' call. It
allocates a new session, reassigns the process, and disconnects the
terminal---all state only the kernel can change. The restriction that
a group leader may not do it exists to keep the group atomic: the OS
refuses to split a group mid-session, so your program must
\emph{fork} to create a child that is not a leader. The detour is the
kernel's safety rule, and daemons are built around it.
\end{osconcept}

\section{Daemons: Processes Without a User}

A \emph{daemon} is a long-running background service: a web server, a
database, a cron scheduler, an inotify watcher. It has no terminal, no
interactive user, no expectation of ever being \code{Ctrl-C}'d from a
shell. It is created to \emph{survive}: it must outlive the login
session that started it, ignore the terminal's death, and keep running
until told to stop.

To become a daemon, a process must undo every link to the user world
it inherited. The standard recipe, in order, is:

\begin{verbatim}
   The daemonization ceremony:

   1. fork()          so the child is NOT a process group leader,
                      and the parent may exit (the child becomes an
                      orphan; init adopts it).

   2. setsid()        the child becomes a new session leader with
                      NO controlling terminal.

   3. fork() AGAIN    now the child is not a session leader either,
                      so it can never reacquire a controlling
                      terminal. The first fork's child exits.

   4. chdir("/")      leave whatever directory we were started in
                      (that mount point may be unmounted later).

   5. umask(022)      set a sane default for files the daemon creates.

   6. close 0,1,2 and re-point all three at /dev/null, so any
      accidental stdout write cannot spam a terminal that no longer
      exists, and stdin reads return EOF instead of blocking.

   The double fork is the part beginners get wrong. Explain it in one
   line: the FIRST fork makes setsid() legal (child is not a group
   leader); the SECOND fork makes the daemon unable to grab a
   controlling terminal again (not a session leader).
\end{verbatim}

The double fork deserves a picture, because it is the most ritualistic
step in all of Unix programming:

\begin{verbatim}
   Why fork twice?

     process A (original)
       |  fork()
       v
     process B   (not a group leader: setsid() is now ALLOWED)
       |  setsid()  -> B becomes session leader
       |  fork()
       v
     process C   (NOT a session leader: can never reacquire a tty)
       |          this is the daemon; B exits immediately
       v
     the daemon lives on, adopted by init, immune to terminals.

   A session leader that later opens a terminal device can be made
   its controlling terminal again. A NON-leader cannot. The second
   fork buys permanent detach.
\end{verbatim}

\begin{lstlisting}[caption={daemonize.c: the whole ceremony in one function.}]
#include <stdio.h>
#include <stdlib.h>
#include <unistd.h>
#include <fcntl.h>
#include <sys/stat.h>
#include <sys/types.h>

/* Re-point descriptor fd at /dev/null. */
static void point_at_devnull(int fd) {
    int nfd;

    nfd = open("/dev/null", O_RDWR);
    if (nfd >= 0) {
        dup2(nfd, fd);
        close(nfd);
    }
}

/* Become a daemon. Returns 0 on success, -1 on failure. */
int daemonize(void) {
    pid_t pid;

    /* 1. First fork: the child may now call setsid(). */
    pid = fork();
    if (pid < 0) {
        return -1;
    }
    if (pid > 0) {
        exit(0);       /* parent's job is done */
    }

    /* 2. New session, no controlling terminal. */
    if (setsid() < 0) {
        return -1;
    }

    /* 3. Second fork: ensure we are not a session leader. */
    pid = fork();
    if (pid < 0) {
        return -1;
    }
    if (pid > 0) {
        exit(0);
    }

    /* 4. Work somewhere safe. */
    if (chdir("/") < 0) {
        return -1;
    }

    /* 5. Default permissions for files we create. */
    umask(022);

    /* 6. Detach all standard descriptors from any terminal. */
    close(0);
    close(1);
    close(2);
    point_at_devnull(0);
    point_at_devnull(1);
    point_at_devnull(2);

    return 0;
}
\end{lstlisting}

After \code{daemonize()} returns, the process has no terminal, no
interactive parent that matters, and no user pressing Ctrl-C. That is
both the point and the danger: nothing will stop it except a signal,
so a well-behaved daemon installs signal handlers (Chapter~6) for
\code{SIGTERM} and, very commonly, treats \code{SIGHUP} as
``reload your configuration'' rather than ``die.''

\begin{osconcept}
Detaching from the user session is entirely an OS ceremony. Forking,
creating sessions, unmounting-proof directories, and re-pointing
descriptors are all kernel services; a user-space process cannot make
itself terminal-independent any other way. (A \emph{mount point} is a
directory where a filesystem is attached: the kernel hangs a disk's
tree at that path. A daemon that keeps working inside a directory
whose filesystem gets unmounted holds a broken view of the world,
which is why the ceremony moves it to \code{/}.) The double-fork rule
exists because the kernel defines who may own a controlling terminal
(only session leaders), so the daemon must manipulate process/session
state
through syscalls to become permanently un-ownable.
\end{osconcept}

\section{Resource Limits}

A daemon runs forever, so it must not grow forever. The kernel enforces
per-process \emph{resource limits}---max open files, max memory, max
CPU time---and your program can read and set them with
\code{getrlimit}/\code{setrlimit}. Each limit has two values: the
\emph{soft} limit (what the process normally gets) and the \emph{hard}
limit (the ceiling the process may raise the soft limit to).

\begin{lstlisting}[caption={Reading and raising the open-file limit.}]
#include <stdio.h>
#include <stdlib.h>
#include <sys/resource.h>

int main(void) {
    struct rlimit rl;
    int r;

    r = getrlimit(RLIMIT_NOFILE, &rl);
    if (r < 0) {
        perror("getrlimit");
        return 1;
    }

    printf("open-file limit: soft = %lu, hard = %lu\n",
           (unsigned long) rl.rlim_cur,
           (unsigned long) rl.rlim_max);

    /* A server wanting more concurrent sockets raises the soft limit
       toward the hard limit. */
    if (rl.rlim_cur < rl.rlim_max) {
        rl.rlim_cur = rl.rlim_max;
        if (setrlimit(RLIMIT_NOFILE, &rl) < 0) {
            perror("setrlimit");
            return 1;
        }
        printf("raised soft limit to %lu\n", (unsigned long) rl.rlim_cur);
    }

    return 0;
}
\end{lstlisting}

\begin{verbatim}
   The limits you will actually need, in one table:

     RLIMIT_NOFILE  how many descriptors you may hold open
                    (a server with 10,000 connections needs a big one)
     RLIMIT_AS      total address space a process may map
     RLIMIT_CORE    size of core dumps (0 = disable them)
     RLIMIT_CPU     seconds of CPU before the kernel kills you

   `ulimit -n` in a shell is the same knob, set by the shell before
   your program even starts. getrlimit() reads it from inside.
\end{verbatim}

\begin{osconcept}
Resource limits are kernel-enforced policy: the kernel counts every
descriptor and every byte of address space against the limit at the
moment it is granted, and refuses to exceed it. A user-space program
cannot ``borrow'' a bit more RAM when the kernel says no. Limits are
the OS's way of containing runaway processes---which is exactly why
daemons both read them and raise them, in full knowledge of what the
kernel will and will not allow.
\end{osconcept}

\section{Practical: A Production-Shaped Daemon}

Let us assemble the chapter---and the signal discipline of Chapter~6---
into a complete daemon: double-forked, detached, signal-handled, with
a \emph{pidfile} (a file holding its pid so \code{stop} scripts can
find it). This is the skeleton that real services are built from.

\begin{lstlisting}[caption={mydaemon.c: the full daemon life cycle.}]
#include <stdio.h>
#include <stdlib.h>
#include <unistd.h>
#include <signal.h>
#include <fcntl.h>
#include <sys/stat.h>
#include <sys/types.h>

static volatile sig_atomic_t g_running = 1;

static void handle_term(int sig) {
    g_running = 0;             /* the only work the handler does */
}

/* Record our pid in a file so `kill $(cat pidfile)` finds us. */
static void write_pidfile(const char *path) {
    int fd;
    char buf[64];
    int n;

    fd = open(path, O_WRONLY | O_CREAT | O_TRUNC, 0644);
    if (fd < 0) {
        return;                /* failure to write the pidfile is not
                                  fatal for the daemon itself */
    }
    n = snprintf(buf, sizeof(buf), "%d\n", getpid());
    write(fd, buf, (size_t) n);
    close(fd);
}

int daemonize(void);           /* defined in Section 9.5 */

int main(void) {
    struct sigaction sa;

    if (daemonize() < 0) {
        perror("daemonize");
        return 1;
    }

    /* Handle stop requests; SIGHUP would normally mean "reload". */
    sa.sa_handler = handle_term;
    sigemptyset(&sa.sa_mask);
    sa.sa_flags = 0;
    sigaction(SIGTERM, &sa, NULL);
    sigaction(SIGINT, &sa, NULL);

    write_pidfile("/tmp/mydaemon.pid");

    printf("daemon started, pid %d\n", getpid());

    while (g_running) {
        /* Real work would go here: accept, read, process, log.
           printf() now goes to /dev/null, so logs must go to a
           file or syslog instead. */
        sleep(1);
    }

    unlink("/tmp/mydaemon.pid");
    printf("daemon stopping\n");
    return 0;
}
\end{lstlisting}

\begin{verbatim}
   Try the full cycle from a shell:

     $ ./mydaemon
     daemon started, pid 12345        (the prompt returns instantly:
                                        the parent exited at fork 1)

     $ ps -ef | grep mydaemon         (it is alive, no tty column)
     $ cat /tmp/mydaemon.pid
     12345
     $ kill -TERM $(cat /tmp/mydaemon.pid)
     $ ps -ef | grep mydaemon         (gone; pidfile unlinked)

   Notice the pidfile made the whole lifecycle scriptable: anyone can
   find the daemon by name, without knowing where it went.
\end{verbatim}

\begin{researchnote}[The evolution of process control]
The session/daemon machinery of this chapter is old, and its design
history is documented by the person who designed much of it. In 1979,
Dennis Ritchie published ``The Evolution of the UNIX Time-sharing
System,'' a retrospective on how the operating system grew: the fork
model, signals, and the shell-based process control we take for
granted were choices made and revised over a decade of use.

\begin{verbatim}
   Ritchie's retrospective, in one line:

       "The first (shell) programs were simple, but the underlying
        process-control mechanisms were general enough to support
        everything that came later."

   Read that as this chapter's thesis: fork, signals, and groups were
   designed as GENERAL primitives, and the shell (and later job
   control and daemons) built enormous behavior on top of a very small
   set of ideas.
\end{verbatim}

It is a short paper and a charming read---part technical, part
biography. It will make the fork/session machinery you just wrote feel
less like arcane ritual and more like the deliberate engineering it is.

\paperref{D.~M.~Ritchie, ``The Evolution of the UNIX Time-sharing
System,'' Language Design and Programming Methodology, 1979}{https://www.bell-labs.com/usr/dmr/www/hist.pdf}
\end{researchnote}

\section{Structures in Bytes}

Most process-control objects live in the kernel and have no stable
byte layout, but two of them have shapes you will actually touch.

\begin{verbatim}
   struct rlimit on 64-bit Linux: 16 bytes
   offset  size  field
   0       8     rlim_cur     the soft limit
   8       8     rlim_max     the hard limit

   a pidfile is plain text:
   "/tmp/mydaemon.pid"  ->  "12345\n"     (pid plus newline)
\end{verbatim}

The kernel's per-process record---its process control block---has a
shape worth knowing even though the offsets change with every version:

\begin{verbatim}
   a minimal view of the per-process record:

   offset  size  field
   0       4     pid
   4       4     ppid
   8       4     pgid       (the process group)
   12      4     sid        (the session)
   16      1     state      (running, stopped, zombie...)
   ...     8     pointer to the open-file table
   ...     8     pointer to the signal dispositions (Chapter 6)
\end{verbatim}

What is stable is the content, not the offsets: a process is a byte
record carrying its ids, its state, and pointers to its tables---the
same facts \code{getpid()}, \code{getpgrp()}, and \code{setsid()}
read and write.

\section{Summary}

\begin{verbatim}
   The process-control model, complete:

   process group   a set of processes signaled as one unit
   session         a set of groups owning one controlling terminal
   foreground group  the group allowed to read the terminal
   setsid          detach: new session, no terminal
   daemon          a process with no user, no tty, long life
   double fork     setsid legality + permanent terminal immunity
   rlimit          kernel-enforced caps on what a process may use
   pidfile         the daemon's way of publishing its pid

   Rules that keep process-control code correct:
     1. background readers get SIGTTIN; foreground groups get SIGINT;
     2. setpgid from BOTH parent and child to beat the race;
     3. a group leader cannot setsid(); fork first;
     4. daemons: chdir("/"), umask(022), close + redirect 0,1,2;
     5. a daemon's only master is a signal: handle SIGTERM/SIGHUP.
\end{verbatim}

\begin{osconcept}
This chapter was the OS as a \emph{sociologist}: grouping processes,
assigning them to terminals, capping their appetite, and adopting
their orphans. Every mechanism was kernel state---groups, sessions,
limits---because only the kernel can enforce rules across processes
that do not trust each other. You are not writing these rules; you are
\emph{negotiating} them through syscalls, and knowing the shape of the
negotiation is what systems programming is.
\end{osconcept}

With process control in hand, we have the full single-machine toolkit:
files, processes, signals, pipes, and daemons. The next chapter turns
to the fast lane of communication---time, timers, and the polling
mechanisms that let one process watch many channels at once without a
thread per channel.

